\section{The accounting}
\label{sec:discussion}
\phantomsection\label{sec:discussion-accounting}\phantomsection\label{sec:discussion-guidance}

Across the full per-intervention accounting (Table~\ref{tab:accounting},
Appendix~\ref{app:accounting}) no single intervention leads on a majority of quantities, and
the one that leads on hold fails the integrity gate, so guidance is keyed to threat model
rather than to a winner. With cooperative users a well-worded context instruction costs
nothing and is the right choice: \emph{discerning} gives the best discrimination of any
prompt-only arm (AUC 0.72 [0.63, 0.81]) at a lower hold (50.3\%), installing on 55.8\% of
topics. Where the deployment can also run the compiled channel, stacking it under that same
prompt keeps the discrimination (0.71) and adds 21 points of hold, the only arm we measured
that is defensible on both axes. Do not reach for a persistence instruction, which raises hold
by producing a model that argues with its own sources. Against an adversarial channel the
compiled arms keep a stance under authority reframes far more often than one-shot instructions
(40--75\% vs.\ 0--4\%), though per-turn re-injection beats both (96\%).%
\phantomsection\label{sec:discussion-cost}
The property that makes a compiled commitment resist that reframe, that it cannot be read as an
instruction, is the same one that prevents an auditor from reading it; the ethics statement
treats that trade as this work's principal dual-use concern.
